\documentclass[11pt]{article}
 
% ---- PACKAGES ----------------


	\usepackage[T1]{fontenc}
	\usepackage[utf8]{inputenc}
	\usepackage[english]{babel}
	\usepackage{graphicx}
	\usepackage{geometry}
	\usepackage{multirow}
	\usepackage{tabularx}			% pour des tableau d'une largeur (totale) définie
	\usepackage{calc}    			% pour calcul d'expressions dans les implément. de longueurs ou compteurs
	\usepackage{dcolumn}
	\usepackage{bm}
	\usepackage[usenames,dvipsnames,table]{xcolor}
	\usepackage{wrapfig}
	\usepackage{fancyhdr}
	\usepackage{array}				% pr tableaux - usuels, e.g. alignemt particulier dans une colonne de type p.
	\usepackage{amsmath}
	\usepackage{amsthm}
	\usepackage{amssymb}
	\usepackage{amstext}
	\usepackage{lastpage}			% To get the last page number
	\usepackage{listings}			% pour utiliser des affichages verbatim améliorés
	\usepackage{authblk}	            % pour les authors dans les abstracts 
%	\usepackage{bbding}  			% pour afficher une icone téléphone
%	\usepackage{wasysym} 			% pour afficher une autre icone téléphone
	\usepackage{marvosym} 			% pour afficher une autre icone téléphone et Email
%	\usepackage{ifsym}   			% pour afficher une icone lettre
	\usepackage{gensymb}				% to display degree symbols \degree
	\usepackage{ifpdf}
	
	% ---- PACKAGE HYPERREF
	
	% ---- CAS PDF
	% *************
	\ifpdf
		% ---- Paramétrage du package hyperref dans ce cas
		\usepackage[%dvips,ps2pdf,hypertex,
		pdftex, 								% Paramétrage de la navigation
		%bookmarks = true, 						% Signets
		%bookmarksnumbered = true, 				% Signets numérotés
		%pdfpagemode = UseNone, 				% None, UseThumbs, UseOutlines, Fullscreen
		%pdfstartview = FitH, 					% FitH, FitV, FitR, FitB, FitBH, FitBV, Fit
		%pdfpagelayout = OneColumn, 			% SinglePage, OneColumn, TwoColumnLeft, TwoColumnRight
		colorlinks = true, 						% Liens en couleur
		linkcolor = blue,
		urlcolor=blue, 
		citecolor = blue,						% Couleur des liens externes
		pdfnewwindow = true						% pour ouvrir les liens vers d'autre pdf dans une nvelle fenêtre
		%pdfborder = {0 0 0} 					% Style de bordure : ici, rien
		]{hyperref}
		%
		%\hypersetup{ 									
		%pdfauthor = {}, 							% Auteurs
		%pdftitle = {}, 							% Titre du document
		%pdfsubject = {}, 							% Sujet
		%pdfkeywords = {}, 							% Mots-clefs
		%pdfcreator = {}, 							% Logiciel qui a crée le document
		%pdfproducer = {} 							% Société avec produit le logiciel
		%plainpages = false}
		%\usepackage{pdfpages}
		%
	
	% ---- CAS DVI
	% *************
	\else
		\usepackage{graphicx}
		\newcommand{\url}[1]{\texttt{#1}}
		\newcommand{\href}[2]{\textcolor{blue}{#2}[1]}
	\fi


% ---- NEW COMMANDS AND ENVS ---------------

\newcommand{\TiOO}{TiO$_2$\,}
\newcommand{\SiOO}{SiO$_2$\,}
\newcommand{\Alumina}{Al$_2$O$_3$\,}

% ---- PATHS ---------------------


% ---- DOCUMENT ------------------

\begin{document}


\begin{center}
\Large{\textbf{Oxide deposition -- Minutes of the Meeting with Ji\v{r}\'{i} Bul\'{i}\v{r}}}
\end{center}
\begin{center}
\large{Thur. 04.08.2016}
\end{center}
\vspace*{0.5cm}


\begin{center}
We met Ji\v{r}\'{i} Bul\'{i}\v{r} from the Department of Analysis of Functional Materials from the Institute of Physics of the AS CR, in Slovanka, Prag. The goal was to know the feasability of the deposition of oxides of controlled thicknesses between 10 and 200~nm, on metals Ti and Al. More precisely, the systems Ti/\TiOO and Al/\Alumina are of interest to study the HSF theory developed by Thibaut J.-Y. Derrien. Duration of the meeting: 1h.
\end{center}

\begin{center}
His prof. website: \texttt{\href{http://www.fzu.cz/~bulir/}{http://www.fzu.cz/$\sim$bulir/}}
\end{center}

\vspace*{5cm}
\tableofcontents

\pagebreak
\section{The Meeting}

\subsection*{What happened?}

\begin{itemize}
\item We exposed what we would like regarding oxides and thickness
\item We exposed the context of the LIPSS and HSFL and the possibility of tuning their spatial period with the oxide layer
\item It was then decided what we can start to do
\item More questions
\item JB showed us his deposition lab
\end{itemize}

\subsection*{What devices do they have?}

They have a lot of devices for depositing layers, making multilayers, controling, in real-time (second time scale), the thickness evolution, characterising the roughness and the properties
\begin{itemize}
\item magnetron sputtering
\item PLD
\item and another one that i forgot (maybe ion beam sputtering, i will ask Inam)
\end{itemize}

\subsection*{Is the guy interested?}

\begin{itemize}
\item It seems that he was interested! He didn't new about LIPSS but he did about plasmons
\item However, he proposed that before going to the whole targets, we will do only 3 for Al/\Alumina to check if it's already working with these 3 points on the SPP curve and if yes we go further.
\end{itemize}



\pagebreak
\section{Practical possibilities}

\subsection*{What thicknesses are possible}

According to him up to 200~nm should not be a problem neither for \TiOO nor \Alumina. He will use magnetron sputtering (if i understood correctly, i will ask Inam also).

Also at the thicknesses we target, it will not further oxidize if in air (which would increase the thickness without control)

\subsection*{What roughness can he get}

\begin{itemize}
\item In his depositions, the roughness strongly depends on the roughness of the initial metallic substrate.
\item He tells that it should be less than 40\,nm and better for thinner films.
\item If he needs to heat the substrate (typically for getting polycrystalline instead of amorphous) the roughness will be badder
\item When growing \Alumina, some kind of independent columns (or kind of pillars) of grown materials are usually seen. this makes the homogeneity not perfect and add some roughness (the top of the columns can be for examples slightly higher than the rest). This make the \Alumina porous... Sorry i understood the concept but really the details in this.
\end{itemize}

\subsection*{Can they characterize the thickness?}

yes, they do with ellipsometry method and can put in the model of calculation roughness. If i understood well their precision is about $\pm5\%$ at low thickness (10-50\,nm) and it increases for higher thickness.

\subsection*{What purity can they get?}

They work hith $10^{-5}\mathrm{Pa}$ vacuum and have very high purity targets, so except the contamination of hydrocarbonates from air before deposition and slightly from the chamber, it should be really high purity.

\subsection*{How to store the sample before irradiation?}

\begin{itemize}
\item basically in some small low vacuum box
\item \Alumina is usally very absorbing water so add dessicative perls in the box
\end{itemize}



\pagebreak
\section{So what should we do?}

We proposed him first to send him a report explaining roughly what we would like to demonstrate and what to be done for this (typically the paper you started on Overleaf git). Then, if he finds it attractive and agrees, we go for a first demonstration with 3 samples, he can also maybe propose something and then if it is encouraging we go further (though he already seems to be ok for the 3 samples).

\subsection*{Send JB the report of what we target}

I told him that we will send him some report on next week. Also that it is not urgent, he goes on holiday soon but don't know when. Note, in this report it could be nice to answer the questions that he asked and was not able to answer:
\begin{itemize}
\item Why these native oxides (\Alumina on top of Al and \TiOO on top of Ti and not others)?
\item Why don't you want roughness in the end?  But now i realize that it should be to be sure about the interference distance?
\item For \Alumina, it crystallizes in different phases with "quite differnt optical properties". Which data did you take for your calculations for \Alumina?
\item For \Alumina, do you want it amorphous, or polycrystalline? 
\end{itemize}

\subsection*{Provide substrates}

We don't have a lot of money remaining (only 1300 EUR) and we will not be able to buy a lot, especially if we go for gold on the other experiment campaign. Luckily, for samples of 1x1cm, he seems okay to give us 3 samples of Al/\Alumina with the desired thickness. He said he has already aluminum. Inam proposed also to bring him the 3-4 Ti samples that he polished. It seems that he said yes, but at the same time he was proposing not to make a lot of samples first. Moreover, although Inam manage to get better roughness than shown in Group meeting, it can still be improved i believe and (see below) roughness of the substrate is critical for the roughness of the oxide surface also. But if we follow these info, then: 
\begin{itemize}
\item JB provides the Al substrate and grow \Alumina on top of it with 3 different thicknesses at the first stage
\item We provide the Ti sbstrate 
\item If your theory also goes for Si/\SiOO, as proposed Inam, we have wafers and can provide Si substrates
\end{itemize}





\end{document}